\documentclass[10pt,a4paper]{amsart}
\usepackage[margin=19mm]{geometry}
\usepackage{amsmath,amssymb,mathtools}
\usepackage[scr=rsfs,cal=euler]{mathalfa}
\usepackage{xfrac}
\usepackage[unicode,hidelinks,bookmarks=false]{hyperref}
\newcommand{\ie}{i.e.\ }
\newcommand{\cf}{cf.\ }
\newcommand{\resp}{respectively\ }
\newcommand{\Subsection}{\subsection}
\providecommand{\nobreakdash}{\nobreak\hbox{-}}
\setlength{\emergencystretch}{2em}
\multlinegap0pt
\usepackage{xcolor}
\usepackage{amssymb}
\usepackage{enumerate}
\usepackage{tikz}
\newtheorem{lemma}{Lemma}
\title{\Large On the structural behavior of images of polynomials}
\author{Tsiu-Kwen Lee, Tran Nam Son}
\date{Source: arXiv:2605.04464v1 (CC BY 4.0)}
\begin{document}
\maketitle

\noindent\emph{Source and licence.} \Large On the structural behavior of images of polynomials. Version arXiv:2605.04464v1 (CC BY 4.0), distributed by arXiv under the Creative Commons Attribution 4.0 International licence, http://creativecommons.org/licenses/by/4.0/, as recorded in the arXiv licence metadata for this exact version and shown on the abstract page https://arxiv.org/abs/2605.04464v1. Full licence notice: \texttt{licenses/arxiv-2605.04464-CC-BY-4.0.txt}.\par\smallskip

\noindent\textit{Editorial scope.} Counted unit: the article's lemma BiSo (source0.tex lines 1200--1205). The article's complete original proof of it is delivered with it (source0.tex lines 1207--1269, one \texttt{proof} environment of 63 lines). No mathematics is written, repaired or re-derived: the statement and the proof are the article's own lines, in the article's own notation, and no retained line is substituted or re-flowed.  No further source-local context is needed: every cross-reference of every retained line resolves inside the two delivered spans.  Nothing was omitted from any retained span, so the package carries no omission record.  No retained line contains a citation, so the wrapper carries no bibliography and no bibliography entry.  No retained line contains a figure environment, an image-inclusion command or drawing code, so no image is delivered and the package declares no asset dependency.  Editorial formatting only, in addition: an AMS \texttt{amsart} preamble that restates the article's own declarations and macros used by the retained lines, namely the environments \texttt{lemma} on separate counters, together with the standard packages those lines need.  The retained environments print in this document as Lemma 1 in order of appearance, whereas the article prints the same items with the numbering of its own full document.  The complete native source of the article as distributed in the arXiv e-print for this exact version is bundled unchanged as \texttt{sources/199863/source0.tex}.
				\begin{lemma}\label{BiSo}
					Let \( D \) be a division ring and let \( n >1  \) be a positive integer. Suppose that \( U \in \mathrm{UT}_n(D) \cup \mathrm{LT}_n(D) \) and \( H = \mathrm{diag}(a_1, \ldots, a_n) \in \mathrm{M}_n(D) \), where \( a_1, \ldots, a_{n-1} \in Z(D) \) and the elements \( a_1, \ldots, a_n \) are pairwise distinct. Then there exist \( P, Q \in \mathrm{GL}_n(D) \) such that
					$
					P^{-1} H U P = H = Q^{-1} U H Q.
					$
				\end{lemma}

				\begin{proof}
					It suffices to consider the case \( A = H U \) with \( U \in \mathrm{UT}_n(D) \), and to argue by induction on \( n \). For \( n = 2 \), let
					$
					A =
					\begin{pmatrix}
						a_1 & y \\
						0 & a_2
					\end{pmatrix}.
					$
					Setting \( x = y(a_1 - a_2)^{-1} \), we obtain
					\[
					\begin{pmatrix}
						1 & x \\
						0 & 1
					\end{pmatrix}
					A
					\begin{pmatrix}
						1 & -x \\
						0 & 1
					\end{pmatrix}
					=
					\begin{pmatrix}
						a_1 & 0 \\
						0 & a_2
					\end{pmatrix}.
					\]
					
					Now assume \( n > 2 \) and that the statement holds for matrices of smaller size. We write
					$
					A =
					\begin{pmatrix}
						a_1 & \alpha \\
						0 & A'
					\end{pmatrix},
					$
					where \( A' \in \mathrm{M}_{n-1}(D) \) is upper triangular with diagonal entries \( a_2, \ldots, a_n \), and \( \alpha \) is a row vector.	Let \( A'' = \mathrm{diag}(a_2, \ldots, a_n) \). By the induction hypothesis, there exists \( P \in \mathrm{GL}_{n-1}(D) \) such that \( P^{-1} A' P = A'' \). Note that \( a_1 I_{n-1} - A'' \in \mathrm{GL}_{n-1}(D) \). 	Set \( x' = -\alpha P (a_1 I_{n-1} - A'')^{-1} \) and define
					\[
					P' =
					\begin{pmatrix}
						1 & 0 \\
						0 & P
					\end{pmatrix}
					\begin{pmatrix}
						1 & x' \\
						0 & I_{n-1}
					\end{pmatrix}.
					\]
					A direct computation shows that
					\[
					P'^{-1} =
					\begin{pmatrix}
						1 & -x' P^{-1} \\
						0 & P^{-1}
					\end{pmatrix}
					\quad \text{and} \quad
					P'^{-1} A P' =
					\begin{pmatrix}
						a_1 & 0 \\
						0 & A''
					\end{pmatrix}.
					\]
					This completes the proof.
				\end{proof}

\end{document}
